\section{Collision rigidity with unregistered launch laws}
\label{sec:unregistered-laws}
The common origin of the homothetic settings can also be removed.
An integrated collision mean first measures the width of the launch
footprint itself. A translation-covariant functional of the complete
collection of centered component supports then separates that footprint
from the obstacles. Neither step requires corresponding obstacles to
be identified across settings. The resulting exact inverse applies to
arbitrary separated convex configurations, before periodicity or a
finite list of component species is assumed.

Let $\mathcal O$ be a nonempty locally finite union of compact convex
bodies with nonempty interiors, diameter at most $D_0$, and mutual
distance at least $d_0$. At setting $i\in\{1,\ldots,m\}$, an unknown
stationary probability density $j_i$ is positive almost everywhere in
the interior of its convex support $A_i$. Suppose the supports are
translated positive homothetic copies of one another. Normalize them as
\begin{equation}\label{eq:registration-model}
 A_i=b_i+\rho_i A_0,\qquad s(A_0)=0,\qquad
 \rho_1=1,\qquad \rho_i>0,
\end{equation}
where $s$ denotes the Steiner point. Thus $A_0=A_1-s(A_1)$ is the
physical first footprint after centering, and $b_i=s(A_i)$. None of
$A_0,b_i,\rho_i$ is supplied. At least two scale factors are distinct.
The densities at different settings need not be rescalings of one
common density. Assume
\begin{equation}\label{eq:registration-separation}
                    2t+\max_i\diam A_i<d_0.
\end{equation}
The observations are the pooled forward and reciprocal reverse means
$F_i,R_i$ at all nominal centers, and $g_i=F_i-R_i$. No regularity
beyond convexity, and no quantitative boundary-mass lower bound, is
required for the exact results of this section.

Write $p_H$ for a laboratory support function and
$p_H^\circ(u)=p_H(u)-s(H)\cdot u$ for its centered version,
$u\in\mathbb S^1$. The centered support functions of all obstacle
components are uniformly bounded and uniformly Lipschitz. The same
is true of their expanded components at any fixed setting.

\subsection{A footprint width without a matching component}
By Theorem~\ref{thm:global-occupation}, $g_i$ determines the occupation
and its complete positive-component collection
\[
               \mathcal P_i=\{P_{i,C}=C+(-A_i):C\in\mathcal C\},
\]
where $\mathcal C$ is the collection of original components.
For any $P\in\mathcal P_i$, choose a bounded domain $E_P$ containing
$P+t\overline B$ and meeting no other component's forward-response
support. Such a domain exists by \eqref{eq:registration-separation}.
Let $\mathcal W(H)$ denote the sum of the two coordinate widths and put
\begin{equation}\label{eq:registration-deficit}
 d_i(P)=\mathcal W(P)-\frac2t\int_{E_P}F_i(x)\dd x.
\end{equation}

\begin{lemma}[The footprint deficit]
\label{lem:registration-deficit}
The quantity in \eqref{eq:registration-deficit} is independent of
$P\in\mathcal P_i$ and of the admissible domain $E_P$. Its common
value is
\begin{equation}\label{eq:registration-width-ratio}
              d_i=\mathcal W(A_i)>0,
              \qquad \rho_i=d_i/d_1.
\end{equation}
In particular, the order and relative sizes of the settings are
determined without any cross-setting component correspondence.
\end{lemma}
\begin{proof}
For $P=C+(-A_i)$ the isolated forward contribution has integral
$t\mathcal W(C)/2$. Indeed a displacement $tv$ adds a strip of area
$t\{p_C(v^\perp)+p_C(-v^\perp)\}$ to the segment sweep; integrating
the launch translation eliminates the probability density, and the
four compass directions count each coordinate width twice. This is
the slice calculation of Lemma~\ref{lem:unk-scale-flux}, which uses
only convexity. Width addition now gives
$d_i(P)=\mathcal W(P)-\mathcal W(C)=\mathcal W(A_i)$.
Translation invariance and homogeneity of width prove the ratio
formula. Nonempty interior makes the denominator positive.
\end{proof}

\subsection{An unlabeled support envelope}
Define the lower centered-support envelope at setting $i$ by
\begin{equation}\label{eq:registration-envelope}
       L_i(u)=\inf_{P\in\mathcal P_i}p_P^\circ(u),
                   \qquad u\in\mathbb S^1.
\end{equation}
This is a finite-valued Lipschitz function, although it need not be
the support function of a convex body and the infimum need not be
attained. It is defined directly from the complete collection; no
ordering, multiplicity convention, or choice of corresponding
components is involved.

\begin{theorem}[A geometric inverse without a supplied origin]
\label{thm:registration-rigidity}
Under \eqref{eq:registration-model}--\eqref{eq:registration-separation},
the exact pooled reciprocal mean pairs determine all scale factors
$\rho_i$, the centered first footprint $A_0$, and the configurations
\begin{equation}\label{eq:registration-output}
                       \mathcal O_i=\mathcal O-b_i.
\end{equation}
They therefore determine the obstacle configuration up to a common
translation and determine the relative footprint translations modulo
its full period group.

Explicitly, choose any $k$ with $d_k\ne d_1$, set
$\rho_i=d_i/d_1$, and put
\begin{equation}\label{eq:registration-support-inverse}
 q(u)=\frac{L_k(u)-L_1(u)}{\rho_k-1}.
\end{equation}
Then $q=p_{-A_0}$ is an actual centered support function. For every
$P\in\mathcal P_i$ the function
\begin{equation}\label{eq:registration-body-inverse}
                         p_P-\rho_iq
\end{equation}
is the support of the corresponding component of $\mathcal O_i$.
Writing $\Pi=\operatorname{Per}(\mathcal O)$, the relative registration
is the exactly determined coset
\begin{equation}\label{eq:registration-coset}
 \{a:\mathcal O_i+a=\mathcal O_1\}=(b_i-b_1)+\Pi.
\end{equation}
No periodicity, finite component count, smoothness, or known homothety
origin is an assumption of this theorem.
\end{theorem}
\begin{proof}
The global occupation inverse gives $\mathcal P_i$, and
Lemma~\ref{lem:registration-deficit} gives the ratios. Steiner centering
commutes with Minkowski addition. With $q_0=p_{-A_0}$ one has
\[
 p_{P_{i,C}}^\circ=p_C^\circ+\rho_iq_0,
 \qquad
 L_i(u)=\inf_Cp_C^\circ(u)+\rho_iq_0(u).
\]
The same infimum occurs at every setting, even when different
components minimize it at different directions. Subtraction proves
\eqref{eq:registration-support-inverse}. Its value is the support of
the true footprint; no smoothness or convexity of the separate
envelopes has been asserted. The identity
$P_{i,C}=(C-b_i)+\rho_i(-A_0)$ proves
\eqref{eq:registration-body-inverse} and hence
\eqref{eq:registration-output}. Formula
\eqref{eq:registration-coset} follows by translating the two recovered
configurations. Their common full period group is also determined
directly from the response by Theorem~\ref{thm:global-response-rigidity}.
\end{proof}

\begin{corollary}[The complete geometric ambiguity]
\label{cor:registration-fiber}
Suppose two admissible experiments have the same pooled reciprocal
mean pairs. Their scale ratios and centered footprints agree, and
there are a vector $h$ and periods $\pi_i\in\Pi$ such that
\begin{equation}\label{eq:registration-gauge}
            \mathcal O'=\mathcal O+h,
              \qquad A_i'=A_i+h+\pi_i.
\end{equation}
Conversely every transformation \eqref{eq:registration-gauge} is
realized with identical collision laws by translating the density at
setting $i$ by $h+\pi_i$. Thus the geometric fiber of the observed
response is exactly this orbit. In a finite nonempty configuration
$\Pi=\{0\}$, so all relative footprint translations are recovered
as vectors, and the only geometric ambiguity is a common translation.
\end{corollary}
\begin{proof}
The theorem gives $\mathcal O-b_i=\mathcal O'-b_i'$ at every setting.
Put $h=b_1'-b_1$. Then $\mathcal O'=\mathcal O+h$, and comparison
at setting $i$ gives $b_i'-b_i-h\in\Pi$. The centered footprints
already agree, proving \eqref{eq:registration-gauge}.

For the converse, couple the new launch displacement to the old one
by $Z_i'=Z_i+h+\pi_i$. Translating both the table and a flight by
$h$ preserves its collision bit, and a further period translation
$\pi_i$ preserves it again. This is a pointwise coupling for forward
and reciprocal reverse preparations, and in fact for every permitted
short direction. It proves equality of the full binary laws, not
merely of their positivity sets. A bounded nonempty configuration
has no nonzero translational period.
\end{proof}

\subsection{Finite periodic reconstruction and registration}
The global envelope is appropriate for the exact full-field inverse.
For a periodic table its role can instead be played by a finite
average. This gives the same finite statistical orders as the
registered experiment.

\begin{lemma}[Cancellation of the periodic patch defect]
\label{lem:registration-period-margin}
Let $Q$ be any fixed convex body and $b$ any vector, and assume the
expanded components $C-b+Q$ remain separated. Their full period group
and primitive component-orbit count are those of $\mathcal O$.
The component comparison distance $d_*$ and the orbitwise patch
defect are unchanged by this common expansion and translation.
Consequently the finite nonperiod margin of
\eqref{eq:patch-margin} applies to period and registration tests on
these expanded components.
\end{lemma}
\begin{proof}
A period permutes complete connected components. Cancelling the same
support function $p_Q$ from the matched pair proves the period
assertion and gives the orbit correspondence. Steiner additivity
also gives, for every pair and every vector $w$,
\[
 d_*(C-b+Q+w,D-b+Q)=d_*(C+w,D).
\]
Take the maximum over one representative of each component orbit and
the infimum over all physical components, in both translation signs.
This defect is independent of the representatives, invariant under
adding a true period to $w$, and unchanged by the displayed expansion.

For completeness, the bounded-patch margin in
\eqref{eq:patch-margin} controls every nonperiod center-difference
class. Given any two components, move their centers separately into
$Q_B$ by presentation-lattice vectors. Their difference changes by
a true period, and its new endpoints lie in the cutoffs of
\eqref{eq:patch-margin}. The defect is unchanged. Moreover $Q_B$
contains representatives of every presentation orbit, so its
maximum is precisely the orbitwise defect. This argument applies
after any common translation, and proves the claimed transfer of
the same positive margin.
\end{proof}

If $n$ is the recovered primitive orbit count, choose one complete
representative per orbit at each setting and define
\begin{equation}\label{eq:registration-orbit-average}
 H_i=\frac1n\sum_{P\bmod\Pi}p_P^\circ.
\end{equation}
The representative choices and their ordering do not matter. Indeed,
\[
 H_i=\frac1n\sum_{C\bmod\Pi}p_C^\circ+\rho_iq,
                 \qquad q=\frac{H_k-H_1}{\rho_k-1}.
\]
Thus finite orbit averages replace the global infima, without
introducing a cross-setting matching problem.

\begin{theorem}[Finite reconstruction of unregistered settings]
\label{thm:registration-finite}
Fix the periodic priors of $\mathcal B_\eta$, a finite number of
settings, uniform $C^{6,\beta}$ centered-footprint and curvature
bounds, and uniform density boundary-mass bounds with exponent
$\gamma\ge0$. Suppose the scales lie in a fixed compact subinterval
of $(0,\infty)$ and at least one satisfies $|\rho_k-1|\ge g_*>0$.
The translations $b_i$ are unknown; a fixed bound on their sizes
suffices to retain a fixed physical launch aperture. Choose a fixed
dyadic compass satisfying the strict separation
\eqref{eq:registration-separation} uniformly over the priors.

A finite controller using only the pooled collision bits recovers
the ratios, centered footprints, a periodic table in the frame
$\mathcal O-b_1$, and representatives of all registration cosets
\eqref{eq:registration-coset}. With probability at least $1-\delta$,
the ratio errors, laboratory $C^2$ errors in this translation frame,
period-basis and free-area errors, and errors of the registration
representatives modulo the true $\Pi$ are at most $C\nu$.
The primitive orbit count and rational period relations are correct.

Let $N_\nu$ count attempted bits, $J_\nu$ count center-labelled
batches, and $S_\nu$ count distinct nominal sites. A later batch at
the same site is counted again in $J_\nu$; each center sampled for
an integrated measurement is one single-attempt batch. Then
\begin{align}
 N_\nu&\le C\nu^{-((\gamma+3/2)s+1)/(s-2)}
           \log(C/(\nu\delta))
           +C\nu^{-2}\log(C/\delta),
                                    \label{eq:registration-attempts}\\
 S_\nu\le J_\nu&\le
       C\nu^{-1/(s-2)}\log(C/\nu)
                  +C\nu^{-2}\log(C/\delta).
                                    \label{eq:registration-batches}
\end{align}
The finest nominal mesh is $O(\nu^{s/(s-2)})$, and the batch list,
including setting labels and repetition counts, has binary
description length at most $C J_\nu\log(C/(\nu\delta))$.
These counts do not assign a manufacturing or metrology cost to the
physical homothety relation or to the realization of nominal controls.
\end{theorem}
\begin{proof}
At each setting run the shrinking-layer rare-collision construction
of Proposition~\ref{prop:shrinking-rare} before footprint subtraction. Its geometric and boundary-mass constants
depend on centered footprint bounds and diameter, not on the
unknown translation. With $h\asymp\nu^{1/(s-2)}$ and $e=c h^s$,
the expanded components are reconstructed with support errors
$O(e)$ in supremum norm and $O(\nu)$ in $C^2$. The fixed physical
aperture includes the unknown translations through their given
coarse bounds. Lemma~\ref{lem:registration-period-margin} permits
period locking at each setting before any calibration. It gives
the common primitive orbit count and the complete representative
lists needed in \eqref{eq:registration-orbit-average}.

Choose any complete component separately at each setting and use
Proposition~\ref{prop:unk-scale-flux-measure} to estimate its
integrated forward mean to error $c\nu$. The proof of that
proposition uses only the isolated expanded component, and hence
applies unchanged to the translated footprint. Form the deficits
\eqref{eq:registration-deficit} and their ratios. Their errors are
$O(\nu)$, since the first footprint width has a fixed positive
lower bound. The separated scale can be selected from these
estimates: a setting maximizing the measured distance from one
has true gap at least $g_*/2$ at sufficiently small error.
The orbit-average formula then estimates $q$ with $C^2$ and
supremum error $O(\nu)$. Subtracting the same $\rho_iq$ from
every component at setting $i$ gives all configurations
$\mathcal O_i$ with the same error. Positive curvature margins
ensure physical output bodies, and the retained finite smoothing
preserves these estimates.

It remains to register the reconstructed frames. Choose a complete
root component of $\mathcal O_i$ and enumerate its center differences
to the finite representative list of $\mathcal O_1$. A true
registration coset has a representative on this list. Test each
candidate by translating all component-orbit representatives and
comparing to the other table modulo its recovered lattice. Only
bounded integer translations can enter the protected test patch,
by the fixed period-basis bounds. An
incorrect candidate differs from a correct one by a nonperiod
center difference in $\mathcal O$. Its defect is at least $\eta$
by Lemma~\ref{lem:registration-period-margin}; the numerical defect
error is $O(\nu)$. Thus the same finite threshold separates the
correct registration coset. Accepted candidates differ by periods
and approximate true coset representatives to $O(\nu)$. No unique
shape, asymmetry, or a distinguished component is required. Use one
reconstructed frame and these registrations to give a common
periodic output. The original period-locking and parallel-area
estimates give the stated discrete, basis, and free-area conclusions.

There are a fixed number of geometric constructions. Each uses
$C h^{-1}\log(C/h)$ center batches, and the shrinking-layer sum of
Proposition~\ref{prop:shrinking-rare} bounds its attempted bits by
$C h^{-1}e^{-(\gamma+3/2)}\log(C/(h\delta))$. A fixed
number of integrated measurements cost
$C\nu^{-2}\log(C/\delta)$ single-attempt batches. Confidence
allocation and fresh preparations give the joint event and both
counts. Orbit averages and registration tests use the finite
reconstructions and require no further collision observations.
The dyadic mesh and the digital description bound are obtained
as in Theorem~\ref{thm:unk-scale-finite}.
\end{proof}

\begin{corollary}[Finite nonperiodic configurations]
\label{cor:registration-finite-cloud}
In place of the periodic prior in
Theorem~\ref{thm:registration-finite}, suppose that $\mathcal O$ has
$1\le n\le n_0$ components, with the fixed diameter, separation,
$C^{6,\beta}$ centered-support and positive curvature bounds.
Assume a known bounded protected aperture contains all expanded
components at every setting and all their forward-response envelopes.
The aperture includes the fixed collars needed for coarse acquisition
and reciprocal inversion. Retain the finite setting count, centered
footprint regularity and curvature bounds, boundary-mass bounds,
compact scale bounds, positive scale gap, coarse bounds on the unknown
translations, and strict compass separation of that theorem.

Without a periodicity or nonperiod-patch-margin assumption, a finite
controller recovers the number $n$, the scale ratios, the centered
footprints, the complete configuration $\mathcal O-b_1$, and the
relative translations $b_i-b_1$ as vectors. With probability at least
$1-\delta$, all scalar errors and the laboratory $C^2$ support errors
in this common translation frame are at most $C\nu$, and $n$ is
correct. The attempted-bit, center-batch and distinct-site bounds
\eqref{eq:registration-attempts}--\eqref{eq:registration-batches},
the nominal precision, and the digital description bound remain
valid, with constants depending also on the complete-aperture and
component-count bounds.
\end{corollary}
\begin{proof}
Coarse acquisition over the known aperture finds every expanded
component. The uniform inballs and separation used in that
construction ensure that each component has retained coarse nodes
and is counted once. The protective collar excludes exterior
fragments, and completeness of the aperture rules out unobserved
components. Thus every setting has the same recovered count $n$,
on the common successful coarse event. This step uses no period
test.

Replace \eqref{eq:registration-orbit-average} by the finite averages
\[
 H_i=\frac1n\sum_{P\in\mathcal P_i}p_P^\circ,
       \qquad z_i=\frac1n\sum_{P\in\mathcal P_i}s(P).
\]
There is one expanded component for each original component, so
\begin{equation}\label{eq:registration-finite-cloud-identities}
 H_i=\frac1n\sum_{C\in\mathcal C}p_C^\circ+\rho_iq,
 \qquad
 z_i=\frac1n\sum_{C\in\mathcal C}s(C)-b_i.
\end{equation}
The deficits again determine the ratios, and any separated scale
gives $q=(H_k-H_1)/(\rho_k-1)$. In addition,
\begin{equation}\label{eq:registration-centroid-registration}
                         b_i-b_1=z_1-z_i.
\end{equation}
These formulas are invariant under independent permutations of the
component lists. Thus neither shape matching nor a discrete
registration decision is required.

Apply the shrinking-layer reconstruction to all components and
the integrated measurement to one arbitrary component per setting,
with the same confidence and accuracy reserves as before. Each
expanded support has $C^2$ error $O(\nu)$ and supremum error
$O(e)$, where $e\asymp\nu^{s/(s-2)}$. Averaging cannot enlarge
these bounds. The scalar deficits have error $O(\nu)$, so the
positive width and scale-gap margins give ratio and footprint
$C^2$ errors $O(\nu)$. Subtraction from the first-setting supports
recovers $\mathcal O-b_1$ with the same error. Steiner points have
error $O(e)$; hence \eqref{eq:registration-centroid-registration}
recovers all relative translations with error $O(e)$, in particular
$O(\nu)$. Finite smoothing preserves these estimates and strict
convexity. The component count is already correct on the coarse
event.

The number of radial searches is bounded by the fixed component
and setting counts times $C\nu^{-1/(s-2)}$. Their shrinking-layer
sum and the fixed number of integrated measurements are exactly
those used to prove \eqref{eq:registration-attempts} and
\eqref{eq:registration-batches}. Complete-component averaging and
the centroid formula require no further observations. The common
dyadic mesh and finite batch encoding therefore give the same
nominal precision and digital description bounds.
\end{proof}
